\chapter{Atualidades e Inteligência Artificial}
\label{cap:atualidades}

\begin{edital}
(1) Tópicos atuais de segurança, transportes, política, economia, sociedade,
educação, saúde, cultura, tecnologia, energia, relações internacionais,
desenvolvimento sustentável e ecologia; (2) Inteligência Artificial:
fundamentos e aplicações --- conceitos de IA, aprendizado de máquina, modelos
generativos e de linguagem (LLMs), ética, governança e privacidade em IA.
\end{edital}
\peso{6 questões, peso 1 --- NOVO em 2026: a disciplina ganhou o bloco de IA.
IA saiu do específico e entrou nas gerais.}

\begin{conceito}[Leia antes: incerteza de FORMATO, não de valor]
Este é o \textbf{único capítulo da apostila sem calibração em provas
anteriores da FGV}. Inteligência Artificial entrou no edital da Dataprev só
em 2026 e não há questão passada da banca para ancorar o estilo. Tudo que
segue sobre \emph{formato} --- se cai como julgamento de afirmativas, como
cenário aplicado, quantas questões são conceituais --- é \textbf{projeção},
não histórico observado como nos demais capítulos.

O que isso muda na preparação: a incerteza é de \textbf{formato, não de
valor}. O conteúdo aqui continua sendo o que a área cobra. E justamente por
ser tema novo, a probabilidade maior é de item \textbf{introdutório/conceitual}
(definições, tipos de aprendizado, o que é LLM, viés, alucinação) do que de
pegadinha sofisticada --- essas a banca só monta depois de anos calibrando o
assunto. Por isso a recomendação é \textbf{dominar bem os fundamentos} desta
parte em vez de tentar antecipar o desenho da questão: fundamento sólido
responde qualquer formato; decorar um formato hipotético, não.
\end{conceito}

\section{Parte 1 --- Atualidades (viés socioambiental)}

Na Dataprev 2024 foram 5 questões (Q31--Q35) e o viés é claríssimo:
\textbf{SOCIOAMBIENTAL}. Sustentabilidade, meio ambiente, ESG, desigualdade,
proteção de dados aparecem quase sempre. A FGV costuma cobrar no formato
\textbf{julgar afirmativas} (V/F, ``está correto o que se afirma em I, II,
III'').

\begin{conceito}[ESG --- as três dimensões]
\textbf{ESG} vem do inglês \emph{Environmental, Social and Governance} ---
em português, \textbf{Ambiental, Social e Governança}. É o conjunto de
critérios usado para avaliar o desempenho \textbf{não financeiro} de uma
organização, hoje presente em relatórios corporativos e em critérios de
investimento:

\begin{itemize}
\item \textbf{E --- Ambiental:} emissões, energia, água, resíduos,
biodiversidade, mudança climática.
\item \textbf{S --- Social:} trabalhadores, diversidade e inclusão, direitos
humanos, comunidades, cadeia de fornecedores.
\item \textbf{G --- Governança:} composição e independência do conselho,
ética, transparência, anticorrupção, prestação de contas.
\end{itemize}

O ``G'' é o mais trocado: a banca oferece ``Gestão'', ``Global'' ou
``Growth'' no lugar de \textbf{Governança}, e às vezes troca o ``S'' de
Social por ``Sustentável'' --- sustentabilidade é o guarda-chuva, não a
letra.
\end{conceito}

Temas concretos vistos nas provas FGV do edital:

\begin{itemize}
\item Fóruns/cúpulas internacionais (G20 no Brasil, presidência rotativa,
``Mundo Justo e Planeta Sustentável'', combate às desigualdades).
\item Impacto ambiental da tecnologia (consumo de energia de data centers,
pegada de carbono do digital --- já caiu na Dataprev 2024, ${\sim}$2\% da
eletricidade mundial no texto-base).
\item Art.\ 225 da Constituição (direito ao meio ambiente equilibrado),
racismo ambiental, justiça socioambiental.
\item Proteção de dados / LGPD aplicada a caso real (uso indevido de dados
de consumidores, publicidade direcionada).
\item No MPU: A3P, Política Nacional sobre Mudança do Clima (Lei
12.187/2009), Protocolo de Quioto, mudanças climáticas (enchentes RS,
incêndios) --- mesma pegada verde.
\end{itemize}

\begin{pegadinha}
Item com \textbf{absoluto} ou termo radical é quase sempre falso: ``todos os
avanços tecnológicos representam INVARIAVELMENTE benefícios para o meio
ambiente'' --- o ``invariavelmente'' já entrega que é falso. Item que mantém
``a lógica de consumo do atual modelo econômico'' e ainda promete ``justiça
social para todos'' é contradição plantada; desenvolvimento sustentável pede
mudar o modelo. Item com dado numérico inflado/inventado (``data centers =
metade do consumo dos ecossistemas digitais'' quando o texto disse
${\sim}$2\% da eletricidade mundial) --- confira o número contra o
texto-base. Item ``bonzinho'' mas falso: destino de dados ``para pesquisas
que democratizam remédios baratos'' soa nobre, mas não é o que a denúncia
trata (uso indevido/comercialização) --- a banca usa a alternativa
moralmente simpática como isca. O item verdadeiro costuma ser o mais
moderado e alinhado ao senso socioambiental.
\end{pegadinha}

\begin{comosair}
Trate como prova de interpretação: a maioria das respostas está no
\textbf{próprio texto-base}. Releia antes de julgar cada item. Marque como
falso todo item com ``sempre, nunca, apenas, invariavelmente,
exclusivamente, todos''. Desconfie de item que contradiz o consenso ESG
(sustentável = mudar consumo, incluir os vulneráveis, reduzir emissões).
Julgue item por item e só então escolha a combinação; se dois itens você tem
certeza, as combinações de resposta já eliminam metade das opções. Acompanhe:
cúpulas do Brasil (G20, COP/clima), LGPD/ANPD em casos noticiados, energia e
IA/data centers, desigualdade.
\end{comosair}

\section{Parte 2 --- Fundamentos de IA (novo, alto retorno)}

\subsection{Conceitos}

\begin{itemize}
\item \textbf{IA:} sistemas que executam tarefas que exigiriam inteligência
humana.
\item \textbf{Machine Learning (ML):} subcampo em que o sistema
\textbf{aprende de dados} e melhora com a experiência, sem ser
explicitamente programado.
\item \textbf{Deep Learning:} ML com \textbf{redes neurais profundas}
(muitas camadas).
\item Relação: IA $\supset$ ML $\supset$ Deep Learning.
\end{itemize}

\subsection{Tipos de aprendizado}

\begin{center}
\begin{tabular}{@{}p{2.6cm}p{2.6cm}p{2.4cm}p{2.7cm}@{}}
\toprule
\textbf{Tipo} & \textbf{Dados} & \textbf{Objetivo} & \textbf{Exemplos} \\
\midrule
Supervisionado & rotulados & prever rótulo & classificação, regressão \\
Não supervisionado & sem rótulo & achar estrutura & clustering (K-Means), associação, PCA \\
Por reforço & recompensa / punição & política ótima & agentes, jogos \\
\bottomrule
\end{tabular}
\end{center}

\begin{pegadinha}
Supervisionado usa dados \textbf{rotulados}; não supervisionado, \textbf{sem
rótulo}. A FGV troca os dois. ``Aprende com os dados e melhora ao longo do
tempo'' = \textbf{redes neurais/ML}, não lógica booleana nem programação
linear.
\end{pegadinha}

\subsection{Ajuste do modelo e métricas de avaliação}

\textbf{Viés (bias) $\times$ variância (variance)} é o trade-off central do
aprendizado supervisionado:

\begin{center}
\begin{tabular}{@{}p{2.6cm}p{4.4cm}p{4.4cm}@{}}
\toprule
& \textbf{Alto viés} & \textbf{Alta variância} \\
\midrule
Modelo & \textbf{simples demais} & \textbf{complexo demais} \\
Sintoma & erra no treino \emph{e} no teste & acerta o treino, erra o teste \\
Nome & \textbf{underfitting} & \textbf{overfitting} \\
Combate & mais atributos, modelo mais expressivo & regularização (L1/L2),
validação cruzada, mais dados \\
\bottomrule
\end{tabular}
\end{center}

Reduzir um tende a aumentar o outro --- por isso é \emph{trade-off}, não uma
escolha em que dá para zerar os dois.

\textbf{Métricas de classificação.} A partir da matriz de confusão
(verdadeiro/falso positivo e negativo):

\begin{itemize}
\item \textbf{Acurácia} = acertos / total. \textbf{Enganosa em base
desbalanceada}: num conjunto com 99\% de e-mails legítimos, o modelo que
chuta ``legítimo'' sempre tem 99\% de acurácia e não detecta \textbf{nenhum}
spam.
\item \textbf{Precisão} = dos que o modelo apontou como positivos, quantos
eram mesmo. Sobe quando o custo do \textbf{falso positivo} é alto (barrar
e-mail legítimo).
\item \textbf{Recall} (revocação/sensibilidade) = dos positivos que existiam,
quantos o modelo pegou. Sobe quando o custo do \textbf{falso negativo} é alto
(deixar passar fraude).
\item \textbf{F1-score} = \textbf{média harmônica} de precisão e recall ---
harmônica, não aritmética, justamente para punir o desequilíbrio entre as
duas.
\item \textbf{ROC/AUC} = capacidade de separar as classes variando o limiar
de decisão.
\end{itemize}

Em \textbf{regressão} (prever um número, não uma classe): \textbf{MAE} (erro
médio absoluto), \textbf{RMSE} e \textbf{R\textsuperscript{2}}. Métrica de
classificação em problema de regressão --- e vice-versa --- é distrator
comum.

\begin{pegadinha}
Inversão de par: dizer que \textbf{alto viés} causa overfitting ou que
\textbf{alta variância} é modelo simples demais. O ancoradouro: viés alto =
o modelo é \emph{burro} (underfit); variância alta = o modelo \emph{decorou}
(overfit). E a regularização combate a \textbf{variância/overfitting}.

Em métricas: chamar F1 de ``média aritmética'' de precisão e recall (é
\textbf{harmônica}); trocar precisão por recall na definição; e tratar
acurácia alta como prova de bom modelo sem olhar o balanceamento da base ---
esse é o cenário que a banca monta com spam, fraude e doença rara.
\end{pegadinha}

\subsection{Modelos generativos e LLMs}

\begin{itemize}
\item \textbf{Modelos generativos:} geram conteúdo novo (texto, imagem)
aprendido da distribuição dos dados.
\item \textbf{LLM (Large Language Model):} modelo de linguagem treinado em
texto massivo (arquitetura \textbf{Transformer}, mecanismo de
\textbf{atenção}); gera/entende linguagem. Ex.: família GPT, Claude, Gemini.
\item \textbf{RAG (Retrieval-Augmented Generation):} combina o LLM com
busca em base externa para respostas mais atuais e fundamentadas.
\item \textbf{Alucinação:} o modelo gera resposta plausível porém
incorreta.
\item \textbf{Prompt, fine-tuning, embeddings} são conceitos recorrentes.
\end{itemize}

\subsection{Ética, governança e privacidade em IA}

\begin{itemize}
\item \textbf{Viés algorítmico} (dados enviesados $\to$ decisão injusta),
transparência/explicabilidade, responsabilização.
\item \textbf{Privacidade:} IA e LGPD (dados pessoais em treino e
inferência).
\item \textbf{Regulação:} ver a subseção seguinte --- é o ponto do bloco que
mais muda até outubro.
\end{itemize}

\subsection{Regulação da IA --- o que está em vigor e o que ainda é projeto}

\begin{conceito}[AI Act europeu: já é lei; PL 2338: ainda é projeto]
A distinção entre os dois \textbf{estados} é o que a banca cobra, e é a mais
fácil de errar por leitura de jornal.

\textbf{União Europeia --- Regulamento (UE) 2024/1689 (``AI Act'').} Não é
proposta: \textbf{está em vigor desde 01/08/2024}, com aplicação
\textbf{escalonada} (as proibições vieram primeiro; as obrigações de sistemas
de alto risco entram depois). Estrutura-se por \textbf{nível de risco}:
\emph{risco inaceitável} (práticas proibidas --- escore social, manipulação),
\emph{alto risco} (obrigações pesadas de conformidade, documentação e
supervisão humana), \emph{risco limitado} (dever de transparência: avisar que
se está falando com uma IA) e \emph{risco mínimo} (livre). É um
\textbf{regulamento}, portanto aplicável diretamente nos Estados-membros, sem
precisar de lei nacional.

\textbf{Brasil --- PL 2338/2023.} \textbf{Ainda é projeto de lei.} Foi
\textbf{aprovado pelo Senado em 10/12/2024} e tramita na \textbf{Câmara dos
Deputados}, onde a votação em plenário foi pautada e adiada mais de uma vez.
Copia o desenho europeu: classificação \textbf{por nível de risco} (excessivo,
alto, e os demais), \textbf{direitos das pessoas afetadas} (informação prévia,
\textbf{explicação} da decisão, \textbf{contestação} e revisão humana), cria o
\textbf{SIA} (Sistema Nacional de Regulação e Governança de IA) sob
coordenação da ANPD, e prevê multa de até \textbf{R\$ 50 milhões por
infração}.
\end{conceito}

\begin{pegadinha}
Três trocas de estado, e todas passam despercebidas. (i) Dizer que o
\textbf{Brasil já tem} marco legal de IA em vigor --- não tem; o que está em
vigor e trata de IA por tabela é a \textbf{LGPD} (art.\ 20). (ii) Dizer que o
\textbf{AI Act ainda é uma proposta} em discussão --- é regulamento vigente
desde agosto de 2024. (iii) Atribuir ao PL 2338 a abordagem ``baseada em
princípios, sem classificação de risco'' --- ele é \textbf{baseado em risco},
como o europeu. Cuidado também com o número: os \textbf{R\$ 50 milhões por
infração} aparecem no PL 2338, na LGPD (art.\ 52) e no ECA Digital --- é o
teto que a banca usa para confundir as três normas entre si.
\end{pegadinha}

\textbf{Atenção à data desta página.} Legislação de IA é o conteúdo mais
volátil da apostila. Confirme o estágio do PL 2338 na semana da prova: se a
Câmara aprovar e houver sanção antes de 11/10/2026, ele deixa de ser
``projeto'' e passa a ser lei nova --- exatamente o tipo de atualização que a
FGV adora cobrar no primeiro concurso seguinte.

\subsection{IA aplicada a negócios e políticas públicas --- como a FGV cobra}

O edital não pede só a teoria (redes neurais, tipos de aprendizado): ele pede
\textbf{fundamentos e aplicações}. Isso significa que a FGV pode enquadrar a
mesma pergunta técnica dentro de um \textbf{cenário} --- um órgão público, um
banco, uma prefeitura --- e cobrar se o candidato reconhece o conceito por
trás da situação descrita. É o mesmo estilo de pegadinha por cenário que já
aparece em Governança (Capítulo \ref{cap:governanca}) e em Segurança
(Capítulo \ref{cap:seguranca}), agora aplicado a IA.

\begin{conceito}[Padrões de cenário que a FGV usa]
\begin{itemize}
\item \textbf{Triagem/priorização de atendimento} (fila de posto de saúde,
concessão de benefício social, análise de currículos): pede se o modelo é
\textbf{classificação supervisionada} (rótulo conhecido: aprovado/negado,
prioritário/não) e se há \textbf{direito a explicação/revisão humana} da
decisão automatizada.
\item \textbf{Detecção de fraude ou anomalia} (glosa de conta pública,
fraude em benefício, transação bancária suspeita): em geral
\textbf{não supervisionado} ou semissupervisionado, porque o padrão fraudulento
muda com o tempo (concept drift) e casos rotulados são raros.
\item \textbf{Chatbot/atendimento ao cidadão} com LLM: cenário testa se você
sabe que a resposta pode \textbf{alucinar} e que, por isso, decisão que afeta
direito do cidadão exige \textbf{validação humana}, não a resposta do modelo
sozinha.
\item \textbf{Concessão de crédito, seguro ou benefício por score
algorítmico}: cenário clássico de \textbf{viés algorítmico} --- se o dado de
treino reflete desigualdade histórica (ex.: bairro, gênero, raça
correlacionados a inadimplência passada), o modelo \textbf{reproduz e
amplifica} essa desigualdade, mesmo sem usar o atributo sensível diretamente
(proxy discriminatório).
\end{itemize}
\end{conceito}

\textbf{Base legal que ancora esses cenários (LGPD, art.\ 20):} o titular tem
direito a \textbf{solicitar revisão} de decisões tomadas unicamente com base
em tratamento automatizado que afetem seus interesses, incluindo decisões
destinadas a definir o seu perfil pessoal, profissional, de consumo e de
crédito. A FGV cruza esse artigo com IA quando o cenário é score de crédito,
RH automatizado ou benefício público.

\begin{pegadinha}[Como a FGV arma a pegadinha em cenários de IA aplicada]
\begin{itemize}
\item O cenário descreve um sistema que \textbf{decide sozinho, sem qualquer
revisão humana}, e a alternativa correta chama isso de boa prática ---
\textbf{descarte}: LGPD garante revisão humana de decisão automatizada
relevante.
\item ``O modelo de IA é neutro porque não usa o atributo raça/gênero
diretamente'' --- \textbf{falso}: variáveis \textbf{proxy} (CEP, escola,
histórico de consumo) podem recriar o viés indiretamente. Ignorar isso é o
distrator clássico de ``neutralidade tecnológica''.
\item Confundir \textbf{explicabilidade} (o modelo consegue justificar a
decisão de forma compreensível a humanos) com \textbf{acurácia} (o modelo
acerta muito) --- um modelo pode ser muito acurado e pouco explicável (caixa-preta), e o edital cobra exatamente essa tensão.
\item Tratar chatbot/LLM em atendimento público como fonte \textbf{definitiva
e infalível} de informação --- ignora o risco de alucinação, que é
justamente o ponto que a banca quer que você saiba mitigar (revisão humana,
RAG com fonte oficial, escopo restrito).
\end{itemize}
\end{pegadinha}

\subsection{Interseção IA $\times$ ESG: energia, governança e regulação}

A FGV gosta de cruzar o bloco de IA com o viés socioambiental da Parte 1
deste capítulo --- é a interseção mais provável do bloco ``Atualidades e IA''
em 2026, porque une os dois assuntos que a banca mais repete.

\begin{conceito}[O que sustenta esse cenário]
\begin{itemize}
\item \textbf{Consumo de energia dos data centers:} a maior fatia da energia
de um data center vai para a \textbf{infraestrutura de TI} (processamento,
${\sim}$metade do consumo) e para o \textbf{resfriamento/ar-condicionado}
(a segunda maior fatia) --- refrigeração é, depois do próprio processamento,
o maior custo energético de treinar/rodar modelos de IA em escala.
\item \textbf{PUE (Power Usage Effectiveness):} métrica-padrão de eficiência
energética de data center (energia total consumida $\div$ energia usada só
pela TI). Regulações recentes (ex.: diretiva europeia) já fixam metas
numéricas de PUE para novos data centers --- quanto mais próximo de 1,0,
mais eficiente.
\item \textbf{Matriz energética:} boa parte da energia elétrica que alimenta
data centers no mundo ainda vem de fontes fósseis --- por isso a expansão
acelerada de IA pressiona diretamente as metas de descarbonização (o mesmo
tema de mudança climática que já caiu em Atualidades).
\item \textbf{Governança:} reguladores discutem exigir que empresas
\textbf{divulguem publicamente} o impacto ambiental (energia, água,
emissões) de produtos e serviços baseados em IA --- transparência como
instrumento de política pública, não só como boa vontade corporativa.
\end{itemize}
\end{conceito}

\begin{pegadinha}
O texto-base cita um número (ex.: ``X\% da energia global'', ``Y milhões de
toneladas de CO$_2$'') e a alternativa \textbf{infla ou distorce esse
número} --- releia o texto-base antes de julgar, nunca confie no número de
memória. Também é comum a alternativa afirmar que ``a infraestrutura de IA
usa MAJORITARIAMENTE energia limpa/renovável'' --- trate como distrator
otimista, salvo se o próprio texto-base afirmar isso explicitamente. E cuidado
com a armadilha inversa: dizer que IA é ``incompatível com sustentabilidade''
também é absoluto demais --- o texto costuma apresentar o dilema (benefício
tecnológico \emph{e} custo ambiental), não um veredito fechado.
\end{pegadinha}

\begin{jacaiu}
\textbf{Em prova real da FGV:} \textbf{supervisionado $\times$ não
supervisionado} (duas questões, uma delas de associação tipo-técnica);
\textbf{viés $\times$ variância} com \textbf{regularização Lasso/L1} e
over/underfitting; \textbf{concept drift $\times$ data drift} num modelo de
\emph{churn} degradando em produção; \textbf{MLOps} num \emph{pipeline}
ponta a ponta de detecção de fraude; \textbf{MAE} para prever tempo de
tramitação (regressão); \textbf{CNN} e camadas de convolução; plataforma
baseada em \textbf{LLM} para análise documental --- \textbf{TJ-RJ} (9
questões de IA). \textbf{Cúpula do G20 no Brasil} (``Mundo Justo e Planeta
Sustentável''); \textbf{impacto ambiental do digital} e consumo de data
centers; \textbf{art.\ 225 da Constituição} e racismo ambiental;
\textbf{LGPD em caso real} (tratamento indevido de dados de consumidores);
\textbf{cidade-esponja} --- \textbf{Dataprev 2024} (Q31--Q35). No MPU, a
mesma pegada verde em outro formato (A3P, Política Nacional sobre Mudança do
Clima, Protocolo de Quioto).

\textbf{No nosso banco} (previsto pelo edital, ainda não visto na amostra de
provas): a hierarquia IA $\supset$ ML $\supset$ Deep Learning; LLM e
Transformer; RAG; alucinação; \textbf{viés algorítmico} em decisão
automatizada (o ``viés'' que caiu no TJ-RJ é o \emph{estatístico} do
trade-off, conceito diferente); \textbf{direito à revisão humana} da LGPD,
art.\ 20 --- a explicabilidade do art.\ 20 caiu, mas no bloco de
\textbf{Legislação} do TJ-RJ, não aqui; a abordagem baseada em risco do AI
Act e o PL 2338/2023; ESG por extenso; PUE e matriz energética de data
center; métricas de classificação.

Rode \texttt{./quiz.py atualidades}.
\end{jacaiu}

\section{Alta probabilidade / pesquisa extra}

\begin{itemize}
\item \textbf{Regularização L1 (Lasso, zera coeficientes) $\times$ L2 (Ridge,
encolhe sem zerar)}, \textbf{concept drift $\times$ data drift} e
\textbf{MLOps} caíram nas provas de IA do TJ-RJ e estão detalhados no
Capítulo \ref{cap:orfaos}, seção de IA/ML --- que também é o capítulo para
onde o quiz manda quem erra questões de \texttt{orfaos}.
\item Acompanhe notícias de IA e regulação de 2026 --- é conteúdo vivo, e
esta disciplina recompensa quem está atualizado na véspera, não quem
estudou em julho.
\end{itemize}
